\quellseite{138}
\section[§ 15. Religious Instruction, 1530–1600]{§ 15. Religious Instruction in the Church Between 1530 and 1600}

E. Sachsse pp. 201 ff. E. Chr. Achelis, pp. 303 ff., 368 ff., 401 ff. G. Schumann u. E. Sperber, \emph{Gesch. d. Religionsunterrichts in der Volksschule}, pp. 35 ff., 1890. G. Langemack, \emph{Historia catechetica} ii, 1733. K. J. Loeschke, \emph{Die religioese Bildung der Jugend u. d. sittl. Zustand d. Schulen im 16. Jahrh.}, 1846. C. Moenckeberg, \emph{Die erste Ausg. v. Luth. Kleinem Katechismus}, 1851, ²1868. J. Brustlein, \emph{Luthers Einfluss auf d. Volksschulwesen u. d. Religionsunterricht}, 1852. F. Ehrenfeuchter, \emph{Zur Geschichte d. Katechismus}, 1857. H. Ernesti, \emph{Zur Orientierung ueber die Katechismusliteratur}, 1859. E. Sachsse, \emph{Zur Gesch. d. ev. Katechismus} in \emph{Halte, was du hast}, 1889. M. A. Gooszen, \emph{De Heidelbergsche Catechismus}, 1890. A. Ernst u. J. Adam, \emph{Katechetische Geschichte des Elsasses}, 1897. *Friedrike Fricke, \emph{Luthers Kleiner Katech. in s. Einwirkung auf die katechetische Literatur d. Reformationsjahrhunderts}, 1898. R. Neumann, \emph{Der ev. Religionsunterricht im Zeitalter d. Reformation}, 1899. Th. Wotschke, \emph{Brenz als Katechet}, 1900. F. Cohrs, \emph{Katechismen u. Katechismusunterricht in PRE³}, 1901. K. Knoke, \emph{Die deutsch-luth. Katechismen i. d. braunschweig-hannoverschen Landen waehrend d. 16. Jahrh.} in \emph{Zeitschr. fuer niedersaechs. Kirchengeschichte}, 1901, 1905. *M. Reu, \emph{Quellen z. Geschichte d. kirchl. Unterrichts im ev. Deutschland zwischen 1530 u. 1600}, 8 vols., 1904-1926. O. Frenzel, \emph{Zur katechetischen Unterweisung i. Zeitalter d. Reformation u. Orthodoxie}, 1915. M. Lauterburg, \emph{Der Heidelberger Katechismus in PRE³}, 1901. A. Lang, \emph{Der Heidelberger Katechismus u. vier verwandte Katechismen, mit einer historisch-theologischen Einleitung} (here more literature concerning the catechisms of the Reformed Church), 1907. Chr. Moufang, \emph{Katholische Katechismen d. 16. Jahrh. i. deutscher Sprache}, 1881. P. Bahlmann, \emph{Deutschlands kath. Katechismen bis z. Ende d. 16. Jahrh.}, 1894. O. Braunsberger, \emph{Entstehung und erste Entwickelung der Katechismen des Petr. Canisius}, 1893. F. X. Thalhofer, \emph{Entwicklung d. katholischen Katechismus in Deutschland}, 1899. A. Richter, \emph{Ev. Kirchenordnungen d. 16. Jahrh.}, 2 vols., 1846. H. Heppe, \emph{Gesch. d. deutschen Volksschulwesens}, 5 vols., 1858-1860. *A. Vormbaum, \emph{Die Schulordnungen d. 16. Jahrh.}, 1860. F. Paulsen, \emph{Gesch. d. gelehrten Unterrichts vom Ausgang d. Mittelalters}
\quellseite{139}
 i, ²1896. *G. Mertz, \emph{Das Schulwesen d. deutschen Reformation i. 16. Jahrh.}, 1902. E. Sehling, \emph{Die evang. Kirchenordnungen d. 16. Jahrh.}, 5 vols., 1903 ff. *Fr. Seefeldt, \emph{Zur Entstehung d. biblischen Geschichts-unterrichts im deutschen Protestantismus}, 1924.

When Luther laid before the Church his Catechism, he was convinced that it met the needs of \emph{oral instruction} of the common people. The very name “Catechism” which he had chosen for this booklet indicates that it was not intended to be anything but a set of oral lessons in literary form. But it was far from his mind to have produced a standard book which must be used in all evangelical churches and schools of Germany. A considerable number of independent catechetical publications made their appearance even after Luther had issued his catechism; and this applies not only to Southern Germany where, e.g., the influential preachers of Strassburg published their catechisms (Butzer in 1534, 1537, and 1543; and Matthaeus Zell in 1535, 1536, and 1537), or to sections largely Reformed, but even to purely Lutheran territory. It is true that most Lutheran catechisms, as those by Moibanus Rhegius, Melanchthon, Lossius, Trotzendorf, Chytraeus, Wigand, were intended for the higher classes of the Latin schools and therefore composed in Latin; but several independent text-books, as those by Chr. Vischer, Hermann Bonnus, Judex, and others, were issued for the primary classes of the German schools and for the instruction of the common people in general. Gradually, however, Luther’s Small Catechism gained predominance not only because Luther was its author but chiefly because it surpassed all others in excellence and because it formed a suitable counterpart to the Heidelberg Catechism (1563) which had welded the various branches of the Reformed Church into a unit. It was only in the upper classes of Latin schools that systematic text-books of Christian doctrine could maintain themselves. Eventually all codes of regulations for churches and schools prescribed the study of Luther’s Catechism
\quellseite{140}
; the Lutheran \emph{Corpora Doctrinae}—these precursors of the Book of Concord of 1580— incorporated it, and after the appearance of the Crypto-Calvinistic \emph{Wittenberg Catechesis} (1571) it found a place among the Symbolic Books of the Church in the Book of Concord. In the schools it was used as First Reader: it contained the alphabet and some reading exercises and offered, in accordance with Luther’s advice, first the catechetical texts and then the texts with Luther’s explanation. In the Latin schools the same was done, but the medium of instruction was, of course, Latin. The advanced pupils, especially in the Latin schools, were given a rather copious analysis and exposition of Luther’s explanation. Among the best of such explanations of Luther’s Catechism in the German language are those by Moerlein (1547), Friedrich (1572), Mathesius (1574), and Bischoff (1599). The most excellent explanation is John Tetelbach’s \emph{Gueldene Kleinod} (Golden Jewel) of the year 1568; here an expository principle of great importance is correctly applied, viz., that an explanation of the Catechism must introduce into, and unfold, the treasure of religious truth contained in the \emph{verbal contents} of Luther’s Catechism. Tetelbach’s booklet was widely used far into the seventeenth century. In the higher classes of the Latin schools the Catechism was taught in Greek; even a polyglot edition (German, Latin, Greek, Hebrew) was issued, and not only once but repeatedly. A German-Latin edition had been prepared soon after the publication of the German Catechism and enjoyed great popularity. Yes, it happened that even in universities it formed the basis for the discussions. It was not long till Luther’s Catechism had spread beyond the boundaries of Germany. Within about 60 years after its appearance it was found in Bohemian, Danish, English (in the form of the Nuremberg Sermons for Children), Finnish (?), Esthonian, French, Greek (three different translations), Hebrew, Icelandic (in
\quellseite{141}
the form of the Nuremberg Sermons), Italian, Lettish, Lithuanian, Latin (four translations), Dutch, Old-Prussian, Slovenian, Spanish (?), Swedish, and Wendish. In but few regions the Catechism of Brenz, in the briefer form of 1535, maintained itself for some time, in Wuerttemberg down to the present day in the form of 1681 which combines with it Luther’s Catechism. Recent specific investigations in this field have disclosed the faithful energy with which the Church in this period everywhere disseminated the Catechism among the young and common people.

The Church, however, did not only provide the needed text-books, but also the institutions required for teaching. Again and again, the home was reminded of its God-given duty of instruction with the result that the inculcation of the Catechism became an important element in family-worship; children and servants would often know the Catechism by memory without, or before, being able to read. In the Latin schools the inculcation and explanation of the Catechism was a regular part of the curriculum, and the school regulations almost without exception contained a specific provision concerning such instruction. In the villages it became an established rule, even beyond Saxony, that the vergers (cf. p. 90) impart religious instruction. Here we have the faint beginnings—not more than that—of the public school. Regular inspections contributed toward preserving and improving these institutions; one code of regulations even required the vergers to give four hours’ instruction daily which was, of course, an excessive demand and for the time being simply could not be enforced. Catechetical instruction was given also in Church: every Sunday the Catechism or one of its Chief Parts was read from the pulpit, special catechetical lessons were conducted on Sunday afternoons, now and then also on week-days, and catechetical sermons were preached during the ember seasons or throughout the
\quellseite{142}
year on Sundays or week-days. All the extant catechetical sermons intended for children are excelled by the above mentioned Nuremberg Sermons for Children; they were translated into several languages (into English in 1548), and in some respects may serve as models even today. When in these institutions the children had learned to understand the Catechism, an examination, called \emph{Katechismus examen} or \emph{Glaubensexamen}, was conducted (as a rule during Lent and usually in the parsonage) which determined whether or not the children might be admitted to Holy Communion. This examination was regarded as the genuine evangelical confirmation. At some places it was preceded by a special course of instruction ranging over several weeks (Schleswig-Holstein 1544, Brandenburg-Ansbach 1564). Gradually public examination before the assembled congregation supplanted the more private examination in the parsonage; it was connected with a simple confession on the part of the children and prayer on the part of the congregation (\emph{exploratio, confessio, oratio}). In this form it was prescribed, for the first time, in 1535 for Liegnitz, and proposed in the “Wittenberg Reformation” of 1545; later it spread also to distant places, but it was by no means universally accepted (Electoral Saxony, e.g., did not introduce it). Chemnitz, in \emph{Examen Concilii Trident.} ii 3, strongly recommended it as the true evangelical counterpart of the Catholic confirmation. In Hesse (1539, 1566, 1574) and in Strassburg (at least since 1543), owing to Butzer’s influence, the imposition of hands was connected with this act and also the vow of obedience to the Church of Christ. Even sacramental character was imputed to confirmation since the reception of the Holy Spirit was identified with the imposition of hands; fortunately, however, this form of confirmation did not spread very far. Further details concerning this point will be found in § 39.

The first communion thus clearly developed as goal of
\quellseite{143}
instruction. It must be admitted, however, that in many cases rote learning of the five Chief Parts, or even of the catechetical texts, was deemed sufficient. At other places the five Chief Parts were supplemented by special questions pertaining to confession and the Lord’s Supper.

Finally, private confession which was practiced almost universally (cf. Augustana xi) enabled the pastor to revive his parishioners’ knowledge of what they had learned in early youth in as much as he examined them in the five Chief Parts of the Catechism; the same was done, here and there, in the examination preliminary to marriage. Now and then circumstances rendered necessary an examination, in private confession, of all adults, but as a general rule to be applied everywhere it could not be justified because it would militate against the Evangelicals’ aim of all educational efforts, viz., the maturity of the Christian congregations.

In the era under discussion also the other tendencies developed that had manifested themselves before and since 1529. In many Latin schools \emph{Biblical History} was taught although it was not yet placed on the same level with the Catechism. Castellio’s excellent \emph{Dialogi Sacri} were edited more than sixty times between 1543 and 1600—although not always for the sake of instruction in \emph{Biblical History}, and eleven editions of Fabricius’ \emph{Historia Sacra} were issued between 1564 and 1600. The same purpose was served by the study of Pericopes which had been introduced in many schools as part of the regular curriculum; in fact, instruction in the pericopes which had been introduced in nearly all Lutheran regions in the 16th century and which was imparted generally, though not exclusively, on Saturdays, was the form of instruction in Bible History in the 16th and 17th centuries. Formerly it was customary among historians to deny that instruction in Bible History was given in this period; but recent research has revealed the error of this view.
\quellseite{144}
More yet, it has been shown that also those children who did not attend Latin schools, frequently studied the pericopes, i.e., Bible History; occasionally the pastors would even explain the pericopes in the churches to children and young people. Furthermore, \emph{collections of Bible Passages} were used in lower as well as higher schools; some of these collections followed the plan of the Catechism, some the order of pericopes. Aside from the feeble beginnings in the twenties, the earliest selections in Latin were prepared by Lossius of Lueneburg and by Trotzendorf in Silesia; Veit Dietrich of Nuremberg composed the first German selection for children (1546). We also find school editions of the Proverbs of Solomon, the Psalms, and Ecclesiasticus. Only by recent research has this wealth of catechetical literature been rescued from oblivion. The young were even taught to sing \emph{hymns} and had to commit some of them to memory; we find hymnals specially designed for children and also catechisms with an appendix containing the chief \emph{hymns}. The diligent and efficient work done between 1530 and 1600 is simply astounding and, in part at least, a model for today.

Faithful work was done also within the \emph{Reformed Church}. As early as 1522, Zwingli instituted a sort of \emph{Kinderlehre} in Zuerich. Here Leo Jud, adopting Luther's arrangement of the Chief Parts, wrote his Catechism in 1534 which was issued in abridged form in 1535. For the church of Geneva Calvin composed one French Catechism in 1537 and another one in 1541 which in the Latin translation of 1545 and in the German translation of 1563 was employed in various districts of Germany. Consisting of 373 questions and answers it was divided into four parts: (1) Faith; (2) the Law; (3) the Lord's Prayer; (4) the Word of God and the Sacraments. It was further divided into 52 sections for the purpose of studying it completely in one year. John a Lasco's Catechism of 1546 (printed in 1551) is exceptionally lucidly arranged
\quellseite{145}
 and in exemplary manner attempts to arouse faith; it contained the following parts: (1) Decalogue; (2) Creed; (3) Baptism, Communion, and Church Discipline; (4) Prayer. All these catechisms were superseded by the Heidelberg Catechism which by electoral rescript was introduced in 1563 in all churches and schools of electoral Palatinate; since the Synod of Dort (1618) it has enjoyed almost symbolical dignity throughout the Reformed Church. Composed by the Heidelberg theologian Ursinus in collaboration with others, notably Olevian, this catechism represents a skillful summary of the previous catechetical efforts in the Reformed Church. It consists of three parts: (1) Man’s Misery (question 1-11); (2) Man’s Redemption (questions 12-85; Creed and Sacraments, including the notorious eightieth question); (3) Man’s Gratitude (questions 85-129; Decalogue and Lord’s Prayer). This famous threefold division was patterned after a catechism by a Lutheran author.\footnote{See M. Reu, \emph{Quellen etc.} i 1, pp. 201–203.} The Heidelberg Catechism is a fair statement of the milder type of Reformed teaching, and here and there rises to heights of linguistic power and beauty, particularly in the famous first question; but its lengthiness and clumsiness and especially the hairsplitting and spinous nature of many of its questions are decided disadvantages. In later editions the proof texts which in the original had been merely indicated, were printed in full and the whole was divided into 52 lessons so as to be covered in the course of 52 Sundays. An abridged edition for children was issued in 1585. This catechism marks the conclusion of catechetical efforts within the Reformed Church in Germany. The needed institutions for the teaching of the catechism were also provided, among them the catechetical examination of old and young on Sundays.

By these educational efforts of the reformatory Church even the Roman Church was induced to make some earnest
\quellseite{146}
endeavors toward providing more general and thorough instruction for the young. A considerable number of catechisms appeared in the sixteenth century, especially since the Jesuits (1540) began to devote themselves to this work. Canisius, the father confessor of Ferdinand I, deserves special mention. In 1556 he issued a small catechism (i.e. a synopsis of his \emph{Summa doctrinae Christianae}) under the title \emph{Summa doctrinae Christianae per quaestiones tradita et ad captum rudiorum accommodata}; soon it was issued also in German (\emph{Der Klain Katechismus} etc.). For two centuries this book was the chief catechism of the Roman Church. In 1559 Canisius issued a catechism for higher schools of learning which despite its bulky size bore the title \emph{Parvus catechismus catholicorum}; it became the standard exposition of the smaller Summa. Even the \emph{Catechismus Romanus} (also called \emph{Catechismus Tridentinus}) though composed by order of the Council of Trent and intended as a manual for priests, or the text books subsequently based on this catechism and designated for exclusive use (1598; pope Clement VIII) did not succeed in dethroning the text books by Canisius.

\quellseite{147}
\section[§ 16. Lutheran Orthodoxism]{§ 16. Catechetical Instruction During the Period of Lutheran Orthodoxism (1600–1676)}

G. Langemack, \emph{Historia catechetica} ii. iii, 1733. 1740. Ph. Schuler, \emph{Gesch. d. katechetischen Religionsunterrichts}, pp. 93 ff., 1802. *F. Ehrenfeuchter, \emph{Zur Geschichte d. Katechismus}, 1857. F. W. Bodemann, \emph{Katechetische Denkmale der ev.-luth. Kirche}, (containing the catechism of Tetelbach 1568, for Nuremberg 1628, of Justus Gesenius 1639, of M. Walther 1653, of Glassius 1640, for Mecklenburg 1711), 1861. *G. v. Zeschwitz ii 1. Schuman-Sperber, pp. 39-49. *E. Sachsse, pp. 186 ff. M. v. Nathusius i, pp. 57 ff. E. Chr. Achelis, pp. 319 ff. J. Steinbeck, pp. 37 ff. A. Ernst u. J. Adam, Katechetische Geschichte des Elsasses, 1897. Buenger, \emph{Entwicklungsgeschichte des lutherischen Katechismusgebrauchs in Hannover}, 1912. *O. Frenzel, \emph{Zur katech. Unterweisung im Zeitalter der Reformation und Orthodoxie}, 1915. *Vormbaum, \emph{Ev. Schulordnungen des 17. Jahrhunderts}, 1863. F. Paulsen, \emph{Geschichte des gelehrten Unterrichts}, 21896. F. Sander, \emph{Geschichte der Volksschule}, besonders in Deutschland in K. A. Schmidt, \emph{Gesch. d. Erz. v 3}, 1902. A. Heubaum, \emph{Gesch. d. deutschen Bildungswesens seit Mitte des 17. Jahrh.}, 1915. A. Matthias, \emph{Gesch. d. deutschen Unterrichts}, 1907. Joh. Arndt; F. Hashagen, \emph{Ausgew. Predigten} in \emph{Predigt der Kirche} xxvi, 1894. G. Hoelscher, \emph{J. Arndt in PRE³}, 1897. W. Koepp, \emph{J. Arndt, e. Untersuchung ueber d. Mystik im Luthertum}, 1912. Joh. Val. Andreae: J. Ph. Gloeckler, \emph{Joh. Val. Andreae, e. Lebensbild}, 1886. J. Bruegel, \emph{Joh. V. Andreae in K. A. Schmidt, Gesch. d. Erz. iii 2}, 1892. G. Hoelscher, \emph{J. V. Andreae in PRE³}, 1896. C. Huellemann, \emph{V. Andreae als Paedagog}, 1884. His \emph{Theophilus}, written 1622 (prior to Comenius!) and printed in 1649, may be found in German in Oehler, J. V. Andreae's \emph{Theophilus}, 1878. His \emph{Evangelische Kinderlehre} upon which the Booklet for Confirmands of Wuerttemberg is based, was published 1621; the third edition (1648) is reprinted in Ph. Schuler, op. cit., pp. 329-352. Wolfgang Ratichius: G. Voigt, \emph{Wolfg. Ratichius, der Vorgaenger d. A. Comenius}, 1894. Joh. Mueller, \emph{Handschriftliche Ratichiana in Paedag. Blaetter}, 1878 ff., esp. 1880, pp. 497 ff. J. Lattmann, \emph{Ratichius und die Ratichianer}, 1898. A. Prall, \emph{Paedag. Schriften d. W. Ratichius}, esp. pp. 61 ff., 1902. H. Barnard, \emph{German Teachers and Educators}, pp. 343-370. R. H. Quick, \emph{Educational Reformers}, ch. ix, 1896. Joh. Amos
\quellseite{148}
Comenius: S. S. Laurie, \emph{A. Comenius, His Life and Educational Works}, 1885. W. S. Monroe, \emph{Comenius}, 1900. *J. Kvasala, \emph{J. A. Comenius, s. Leben u. s. Schriften}, 1892. J. Kvasala, \emph{Die paedag. Reform d. Comenius in Deutschland bis z. Ausg. d. 17. Jahrh.}, 1904. F. Pappenheim, \emph{J. A. Comenius}, 2 vols., 1902. 1905. \emph{Monatshefte d. Comeniusgesellschaft}, 1892 ff. His \emph{Didactica Magna} has been translated into German by Lion, ⁵1904, into English by Keatings, 1896; \emph{Schola ludus} ed. by Boetticher, ²1907; \emph{Mutter Schul} ed. by Lion, ²1907, in English: W. S. Monroe, \emph{The School of Infancy}, 1896; \emph{Orbis pictus} ed. by C. W. Bardeen, 1891. Joh. Kromayer: L. Weniger, \emph{Ratichius, Kromayer u. der Neue Methodus in Weimar} in Zeitschr. f. thuering. Geschichte, 1896. L. Weniger, \emph{J. Kromayers Schulordnungen von 1614 u. 1617}, ibid., 1901. Chr. Waas, \emph{Die Generalvisitation . . . von 1641 bis 1647}, ibid., esp. pp. 245 ff., 1909. Sigm. Evenius: P. Stoetzner, \emph{Sign. Evenius, e. Beitrag zur Gesch. d. Ratichianismus}, 1895. G. Schmid, \emph{S. Evenius in Monatshefte d. Comeniusgesellsch.}, 1895. K. Knoke, \emph{E. christl. gottselige Bilderschule} in Katech. Zeitschrift, 1907. Cf. also Lattmann, \emph{Ratichius u. d. Ratichianer}, 1898. Ernest the Pious, Duke of Gotha: G. Kreyenberg, \emph{Ernst d. Fromme, ein Lebens- und Kulturbild des 17. Jahrh.}, 1890. H. Petrich, \emph{Herzog E. d. Fr., s. Leben und Wirken}, 1901. W. Boehne, \emph{Die paedagogischen Bestrebungen Ernsts d. Fr. v. G.}, 1888. H. Schroedel u. H. Moeller, \emph{E. d. Fr., ein Paedagog unter d. Fuersten}, 1901. M. Ehr, \emph{Beitr. z. Kirchen- und Schulverfassung d. Herzogtums Gotha bis z. Tod Ernst des Fr.}, 1891. *Braem, \emph{Der gothaische Schulmethodus}. \emph{E. kritische Untersuchung ueber die ersten Spuren d. Pietismus i. d. Paedagogik d. 17. Jahrh.}, 1897. The Methodus of 1642 has been edited with an excellent critical and historical introduction by Johannes Mueller in \emph{Sammlung selten gewordener paedagogischer Schriften frueherer Zeiten}, No. 10, 1883; it bears the title: “Herzog Ernst des Frommen Spezial und sonderbarer Bericht, wie die Knaben und Maegdelein auf den Dorfschaften und in den Staedten im Fuerstentum Gotha kurz und nuetzlich unterrichtet werden koennen und sollen”. The Methodus of 1672 has been edited by A. Prall, \emph{Der Schulmethodus des Herzogs Ernst v. Gotha}, ²1912. S. Glassius’ \emph{Kurzer Begriff der christlichen Lehre} is reprinted by Ehrenfeuchter, op. cit., Supplem. pp. 41-62. Michael Walther: Ehrenfeuchter, op. cit., pp. 82 ff. R. Buenger, op. cit. His Catechism is reprinted by Ehrenfeuchter, op. cit., supplem. pp. 89-126, and by Bodemann. Justus Gesenius: Ehrenfeuchter, op. cit., pp. 79 ff. G. v. Zezschwitz ii 2,1, pp. 93 ff., 1874. Buenger, op. cit. E. Bratke, \emph{Just. Gesenius, sein Leben u. s. Einfluss auf d. Hannoversche Landeskirche}, esp. pp. 44 ff.; 82 ff.; 146 ff., 1883.
\quellseite{149}
C. Bertheau, Justus Gesenius in PRE³, 1899. K. Kayser, Die General-visitation des Just. Gesenius in Zeitschr. fuer niedersaechs. Kirchen-gesch., 1906. His Catechism is reprinted by Bodemann, op. cit., and in parts by Ehrenfeuchter, op. cit., supplem. pp. 62 ff.

When in opposition to its obscurations and perversions, Lutheran doctrine had been confessionally fixed in the Book of Concord of 1580 or more specifically in the Formula of Concord, the temptation was very enticing to work all this doctrinal material into the catechetical text books though they were not designed for theologians but for children. Generally speaking, the Church succumbed to this temptation. As a result, the doctrinal element was so preponderant in the catechisms of the time that the religious and practical elements were often completely overshadowed. An abhorrent example of this kind is found in Conr. Dietrich’s Institutiones Catecheticae of 1613. In itself it is an excellent book, helpful alike to student and pastor; but the author’s method of dragging into it every technical term of dogmatics and history of dogma has had baneful consequences; even today we are not quite free from this injurious practice. In catechetical instruction the catechist found a welcome opportunity to review his dogmatics. This proceeding may have been salutary for himself, but the children were given stones rather than bread. Even the epitome of Dietrich’s Institutiones designed especially for school use and employed in some regions to our day, is far too doctrinaire. The same criticism applies to the catechetical sermons of this period. The explanation of the Fifth Chief Part in Aegidius Hunnius’ catechetical sermons would be a highly creditable university lecture; others even surpass Hunnius in this respect. Catechetical sermons like those of John Arndt (1620) are rarely found in this era. Noble men were grieved by these conditions; it was especially Valentin Andreae, his soul on fire for the welfare of the Church, who made bitter complaint of the parrot-like recitation of a catechism and an explanation
\quellseite{150}
which were not understood, and severely scourged such practice in his Theophilus (1622). To improve these deplorable conditions he published his Kinderlehre (1621). Compare also § 35.

Then came the Thirty Years’ War with its woeful, its unspeakable devastation of church and school. The educational institutions that had developed in the course of time were dangerously undermined and some of them disappeared entirely, for instance the preparation of the children for their first communion. But God’s dire visitation through war once more filled the hearts with a desire to care for the country’s youth so shockingly neglected during the incessant tempest of battle. Thus a revival of education came to pass in the midst of the turmoil of war. As early as 1628 Nuremberg issued the equally instructive and edifying \emph{Kinderlehrbuechlein} (Booklet for the instruction of children); it must be admitted, however, that this booklet promoted the one-sided evaluation of the Decalogue as a mere mirror of sin by concluding the explanation of each commandment with the question “Have you kept this commandment?” Strassburg, likewise, distinguished itself by taking good care of its youth; especially J. Schmidt and J. K. Dannhauer, professors at the university, took a warm interest in this matter. Stirred by their example, faithful Justus Gesenius of Hanover devoted himself particularly to the improvement of education in his country. In his \emph{Catechetical Questions} which in the issue of 1639 exerted far reaching influence, he induces his pupils to think by changing the form of every question two or three times, and he lays stress upon the application of the truth to life. Also Michael Walther of Celle helped to raise catechetical instruction to a higher level. His explanation of the Catechism which appeared first in 1653 and was reissued in 1861 and since then frequently used in America (the so-called \emph{Stohlmann Catechism}), excels that of Gesenius in clearness
\quellseite{151}
 of definitions and exactness of method, but introduces too much dogmatical material.

In the Thuringian region it was especially general-super\-intendent Kromayer who made earnest efforts at improving religious instruction. Still better, more systematic, and more permanent work was done by the excellent duke Ernest the Pious of Gotha (1601-1675) and the splendid educators Evenius and Reyher. They succeeded in putting into practice the creative thoughts of the eminent pedagogue Ratichius († 1635) and the still greater Amos Comenius (1592-1671) and built up model schools throughout the region. The school regulations of 1642, and subsequent extensions, stipulated that in the villages all boys and girls from five to twelve years of age must attend school; the munificent duke furnished the needed equipment, in part even the school books; he insisted on the employment of efficient teachers; and he was the first to incorporate (1642, resp. 1662) scientific branches, as arithmetic and natural history, in the curriculum of the village school which formerly had included only religion, reading, writing, and singing. Thus in every respect he laid the foundation of the Christian common school of today. As regards religious instruction, Comenius had demanded that all children should learn by heart those Psalms and hymns which were most frequently sung in their respective localities, and that they should be able accurately to recite not only the Catechism but also the Bible stories and the chief passages of Scripture. In Gotha these requirements were met by teaching the children the following subjects: (1) the text of the Catechism; (2) Luther’s explanation; (3) the \emph{Kurze Begriff} by Sal. Glasius (an explanation of Luther’s Catechism); (4) 160 Bible passages; (5) Psalms, hymns, and prayers. In 1636 Evenius (?) or possible Brunchorst had published a sort of introduction into Biblical History under the title “Christian, Godly Instruction through
\quellseite{152}
Pictures, i.e., Instruction of the Younger Children in Godliness through Biblical Pictures,” but it was never employed in the schools; a little later, however, books of Bible stories were issued which at least the older children were (6) expected to read (1662) whereby Bible reading was introduced into the regular work of the school. Subsequently the children were given (7) a “Gospel Book” containing the pericopes taken from the Gospels. These lessons were read in school and explained to the children in order that they might intelligently appreciate the sermons preached on these texts. Finally (8), they were examined at school in regard to the chief contents of the sermon. During the long reign of Ernest, these regulations took root in his domain, and the Christian public schools found today in the land of the Reformation owe their origin, directly or indirectly, to his institutions.

Even in Gotha, this model country, instruction in the Catechism occupied the foreground rather exclusively; Bible instruction, however, was at least part of the curriculum. Elsewhere, too, selections of Bible passages and collections of pericopes were used in the instruction of children; and that instruction in Biblical History was not totally ignored becomes evident from the repeated publication in this era of Castellio’s \emph{Dialogi Sacri} and Fabricius’ \emph{Historia Sacra} (cf. p. 121). The above mentioned Gesenius issued in 1656 a Biblical reader for home and school under the title: “Bible Stories of the Old and New Testaments, Chronologically Arranged for the Benefit of the Young and Unlearned, Divided into Two Parts each with 54 Lessons.” In the preface Gesenius emphasizes the fact that the child’s mind craves stories, that without knowledge of Bible History neither catechetical instruction nor sermons can be properly understood, and that chronological knowledge of the stories of the Bible is indispensable if one would comprehend the progress in the economy of salvation.

\quellseite{153}
Of the better catechisms of this period we mention those of Altenburg, Quedlinburg, Frankfort-on-Oder, and Danzig. Some of these offer also selections of Bible passages and a number of Psalms and hymns. Rector Maukisch of Danzig suggested ways of reforming the method of catechetical instruction (1653), and Theo. Grossgebauer wrote on the betterment of education in general (1661).
